Statistical disclosure control protects sensitive data by restricting what leaves a trusted research environment (TRE). The same restrictions can remove the evidence a reviewer needs to recognize a wrong analysis. We measure this tension as local language models begin to do analytical work inside TREs. The first study asks whether a model can review analyses under increasingly strict disclosure rules: eight kinds of analytic error, each planted in a report and paired with a correct report built from the same tables, are released at five levels of disclosure control and shown to reviewers. The second study evaluates models as analysts before admission: \N{harvest.tasks.total} analysis tasks whose data invite each error are given to open-weight models, their write-ups are judged against hidden answer keys, and other models check those write-ups. Disclosure rules alone detected no analytic error. Two hosted reviewers detected every planted error when nothing was withheld, and lost detection only under the strictest policy and mainly for the errors that need a count. Much of their detection rested on the wording of the interpretation. For model analysts, a prompt that makes the model record the study design before it writes, with a short thinking budget, reduced errors across three Qwen3.5 models, and a larger model added no consistent benefit within that group. Qwen3.5 27B followed by a gpt-oss 20B checker met the study's release threshold on both sealed task sets under every judge. Checkers, however, separated analysts' real errors far less well than planted ones. In the companion gateway, an open-weight planner turned requests into the intended typed query on almost every answerable request once its prompt stated what each view holds, but neither planner tested refused what the query language cannot express: both answered a different, permitted question, so refusal must come from the gateway. These results argue for evaluating disclosure protection, analytical correctness and review performance separately. A companion paper tests the same models for bias and deception. The results are a machine baseline: human review and adjudication are still to come.
